\scanpage{121}
\mspage{59}
% L01
\[
\boxed{\sigma^{*}(F_D)\xrightarrow{\varphi}\rho^{*}(F_U)}
\]
% L02
Or, \deleted{\unsure{cette donnée si $F_U$ est un}} \inserted{si on se borne aux} $F_U$ tels que $F_U$ soit
% L03
lui-= ess$^{t}$. loc$^{t}$ constant, alors les données
% L04
$F_U,F_D,\varphi$ ne font intervenir que des
% L05
faisceaux ess$^{t}$ loc. constants, i.e. des
% L06
topos fondamentaux (multigaloisiens)
% L07
associés aux schémas envisagés, \deleted{\unsure{savoir}} donc
% L08
se décrivent entièrement en termes des
% L09
diagrammes
\diagram{\begin{tikzcd}[ampersand replacement=\&]
 \& B_{\pi_1,D^{*}} \arrow[dl,"{\deleted{\illegible{1}}\,\rho}"'] \arrow[dr,no head,"\sigma"] \& \\
 B_{\pi_1,U} \& \& B_{\pi_1,D}
\end{tikzcd}}
\margin{(*)}
% L10
-- d'ailleurs le cas où $\varphi$ est un \uline{iso} correspond
% L11
justement au cas des faisceaux loc. ess.
% L12
constants sur $X$, i.e. de faisceaux sur
% L13
$B_{\pi_1,X}$ -- lequel topos fondamental
% L14
apparaît donc comme somme
% L15
amalgamée du diagramme de topos
% L16
précédent, s'insérant dans le carré
\diagram{\begin{tikzcd}[ampersand replacement=\&]
 \& B_{\pi_1,D^{*}} \arrow[dl] \arrow[dr] \& \\
 B_{\pi_1,U} \arrow[dr] \& \& B_{\pi_1,D} \arrow[dl] \\
 \& B_{\pi_1,X} \&
\end{tikzcd}}

\scanpage{122}
% L01
Supposons pour simplifier $U$ connexe,
% L02
choisissons un rev. universel \inserted{$\widetilde U$} de $U$, et un
% L03
rev. universel \inserted{$\widetilde D_i^{*}$} de chaque $D_i^{*}$ -- d'où un
% L04
rev. universel \inserted{$\widetilde D_i$} de $D_i$ -- et des isom.
% L05
(= ``classes de chemins'') entre $\rho_{!}(\widetilde D_i^{*})$ et $\widetilde U$
% L06
-- donc le diagramme des topos \unsure{(*)},
% L07
-- ou des groupoïdes fondamentaux --
% L08
s'explicite en termes d'un diagramme
% L09
\diagram{\begin{tikzcd}[ampersand replacement=\&,column sep=large]
 \& {\begin{gathered}
 (\operatorname{Gal}(\overline K_i/K_i))\\[-2pt]
 \Vert\\[-2pt]
 \mathcal E_i
 \end{gathered}} \arrow[dl,"\unsure{\rho_i|\varphi_i}"'] \arrow[dr,"\sigma_i"] \& \\
 {\begin{gathered}
 \mathcal E\\[-2pt]
 \Vert\\[-2pt]
 (\pi_1(U,\widetilde U)=\operatorname{Aut}_U(\widetilde U))
 \end{gathered}} \& \&
 {\begin{gathered}
 \Gamma_i\;\illegible{1}\\[-2pt]
 \Vert\\[-2pt]
 (\operatorname{Gal}(\overline{k(s)}/k(s)))
 \end{gathered}}
\end{tikzcd}}
% L10
et la donnée d'un $F$ comme envisagé sur
% L11
$X$ revient à la donnée d'un système
% L12
$(E_U,(E_i)_{i\in I},(\varphi_i)_{i\in I})$, où $E_U$ est un $\mathcal E$-ens,
% L13
$E_i$ un $\Gamma_i$-ens, $(\forall i\in I)$, et $\varphi_i$ un
% L14
$\mathcal E_i$-hom de $E_i$ dans $E_U$ (i.e. un
% L15
un $\Gamma_i$-hom
\[
\varphi_i:\ E_i\longrightarrow E_U^{\Pi_i},
\]

\scanpage{123}
\mspage{60}
% L01
où \deleted{\unsure{$\Pi_i$ est}} $\Pi_i=\unsure{\operatorname{Ker}(\sigma_i)}$ s'insère dans la suite exacte
% L02
\[
1\longrightarrow\Pi_i\longrightarrow\mathcal E_i\longrightarrow\Gamma_i\longrightarrow1
\]

% L03
Le cas où les $E_i\longrightarrow E_U$ sont des iso,
% L04
i.e. les $E_i$ \unsure{s'identifient|s'identifiant} \deleted{\illegible{1}} tous à $E_U$,
% L05
avec action triviale \unsure{de|des} $\Pi_i$ sur $E_U$, correspond
% L06
au cas où $F$ est ess. loc. \unsure{ct|cst} sur $X$,
% L07
d'où aussitôt
% L08
\[
\pi_1(X)=\mathcal E=\pi_1(U)\ \big/\
\substack{\text{ss-groupe \unsure{invariant} engendré}\\
\text{par les }\rho_i(\Pi_i)}.
\]

% L09
Notre propos est celui d'un calcul galoisien
% L10
(si possible) de la cohomologie de $X$ pour
% L11
les $F$ envisagés.
% L12
Soit $B_{X,U}$ le topos \inserted{dont les faisceaux sont} des triples
% L13
$(F_U,F_D,\varphi)$ comme dessus, qui s'envoie
% L14
donc dans le topos $B_{\pi_1,X}$, correspondant
% L15
aux couples pour lesquels $\varphi$ est iso,
% L16
et reçoit le topos $X_{\mathrm{ét}}$ :

\scanpage{124}
% L01
\diagram{\begin{tikzcd}[ampersand replacement=\&]
X_{\mathrm{ét}} \arrow[r,"f"] \& B_{X,U} \arrow[r,"g"] \& B_{\pi_1,X}
\end{tikzcd}}
% L02
On se pose la question

% L03
a) Calcul ``explicite'', en termes de cohomologie
% L04
des groupes profinis, de la cohomologie
% L05
dans le topos abracadabra $B_{X,U}$. (NB.
% L06
\deleted{\unsure{b) le calcul de}} la cohomologie \deleted{\illegible{1}} de
% L07
$B_{\pi_1,X}$ n'est autre que la cohomologie
% L08
galoisienne profinie de $\pi_1(X)=\mathcal E/\ldots$
% L09
calculé plus haut\ldots).

% L10
b) Vérifier si pour un faisceau de
% L11
torsion $F$ sur $B_{X,U}$ l'hom. can.
% L12
\[
H^{*}(B_{X,U},F)\longrightarrow
H^{*}(X_{\mathrm{ét}},f^{*}F)
\]
% L13
est un iso.

% L14
NB Jusqu'à maintenant, l'hyp.
% L15
$\dim X=1$ n'a pas servi, ni l'hyp. noethérien
% L16
-- seulement le fait que $S$ soit partie fermée
% L17
discrète, $X$ connexe\ldots. Pour les calculs qui

\scanpage{125}
\mspage{61}
% L01
suivent, correspondants du \deleted{\illegible{1}} cas où les
% L02
$\unsure{D_s|O_s}$ sont des \unsure{jauges} à corps résiduel $k(s)$
% L03
de car. $0$, on va supposer que
% L04
\[
\Pi_i\simeq\widehat{\mathbb Z}\qquad\text{(non canoniquement)}
\]
% L05
-- en fait, par la théorie Kummérienne on
% L06
a un système de systèmes locaux de Tate
% L07
$\unsure{T_{\mathcal E}}$ \inserted{(sur $\mathcal E$)}, $T_i$ \inserted{(sur $\Gamma_i$)}, et des \inserted{$\mathcal E_i$-}isom. $T_{\mathcal E}\simeq T_i$, \deleted{et \illegible{1}}
% L08
(i.e. un $T_X$ sur lequel opère $\unsure{\widetilde{\mathcal E}|\overline{\mathcal E}}=\mathcal E/\ldots$)
% L09
et on aura un iso. can. Kummérien
% L10
\[
T_i\xrightarrow[\sim]{\unsure{\kappa_i|k_i}}\Pi_i.
\]

% L11
On aimerait pouvoir paraphraser,
% L12
sur $B_{X,U}$, la suite exacte de cohomologie
% L13
bien connue
% L14
\[
\longrightarrow H^i(X,F)\longrightarrow H^i(U,F)
\longrightarrow H^{i+1}_S(X,F)\longrightarrow\cdots
\]
% L15
relative à l'ouvert $U$ du topos $X$, et son
% L16
``complémentaire'' fermé $S$. Or, \deleted{\illegible{2}} \inserted{tout comme}
% L17
$X$ s'insère dans un diagramme de topos
% L18
\diagram{\begin{tikzcd}[ampersand replacement=\&]
 \& D^{*}_{\mathrm{ét}} \arrow[dl] \arrow[dr] \& \& \\
 U_{\mathrm{ét}} \arrow[dr,no head] \& \& D_{\mathrm{ét}} \arrow[dl] \& S_{\mathrm{ét}} \arrow[l,"\unsure{\xi}"'] \\
 \& X_{\mathrm{ét}} \& \&
\end{tikzcd}},

\scanpage{126}
% L01
de = $B_{X,U}$ s'insère dans
% L02
\diagram{\begin{tikzcd}[ampersand replacement=\&]
 \& B_{\pi_1,D^{*}} \arrow[dl] \arrow[dr] \& \& \\
 B_{\pi_1,U} \arrow[dr] \& \& B_{\pi_1,D} \arrow[dl] \& B_{\pi_1,S} \arrow[l,"\approx","\unsure{\xi_{\mathrm{ét}}}"'] \\
 \& B_{X,U} \& \&
\end{tikzcd}}

% L03
[\uline{NB} les $\varprojlim$ finies et $\varinjlim$ quelc.
% L04
dans $B_{X,U}$ i.e. pour les $(F_U,F_D,\varphi)$, se
% L05
calculent ``terme à terme''] Je dis
% L06
que $B_{\pi_1,U}\longrightarrow B_{X,U}$ s'identifie à un
% L07
morphisme d'induction, relatif à l'
% L08
objet \deleted{\illegible{1}} \unsure{(noté encore} $U$ par abus de notation)
% L09
de $B_{X,U}$ défini par
% L10
$F_U={}$faisceau final $e_U$, $F_D={}$faisceau initial $\varnothing_D$
% L11
(et $\varphi$ \unsure{est|étant} alors fixé !). En effet, les
% L12
objets \inserted{de $B_{X,U}$} au dessus de $U$ s'identifient aux
% L13
$(F_U,F_D,\varphi)$ avec $F_D=\varnothing_D$ (ce que fixe $\varphi$ !) donc
% L14
ils forment une cat. équivalente à celle des
% L15
$F_U$, i.e. à $B_{\pi_1,U}$. On voit que le topos
% L16
résiduel (s'identifiant \deleted{aux} \inserted{à la cat. des} $F\deleted{_{\unsure{U}},\unsure{E|e}_{\unsure{i}}}=(F_U,F_D,\varphi)$ tels

\scanpage{127}
\mspage{62}
% L01
que $F_U=e_U$ faisceau final, s'identifie de = à la
% L02
catégorie \inserted{$B_{\pi_1,D}$} des $F_D$ \deleted{\illegible{2}} -- le foncteur
% L03
canonique ``image inverse sur $B_{\pi_1,D}$ de l'image
% L04
directe de $F$ sur $B_{\pi_1,U}$'' n'étant autre
% L05
$\sigma_*\rho^*$, de sorte qu'on retrouve la description
% L06
\deleted{\illegible{1}} typique d'Artin d'un topos déduit
% L07
par ``recollement'' d'un \deleted{\unsure{objet}} ouvert et du
% L08
fermé complémentaire.
% L09
On trouve donc une suite exacte
% L10
\[
\longrightarrow H^i(B_{X,U},F)\longrightarrow
H^i(B_{\pi_1,U},F_U)\longrightarrow
H^{i+1}_{S(\text{ou }D)}(B_{X,U},F)\longrightarrow\cdots
\]
% L11
s'envoyant dans la suite exacte analogue
% L12
relative à $X_{\mathrm{ét}},U_{\mathrm{ét}},S_{\mathrm{ét}}$. \deleted{\unsure{Pour} \illegible{2}} \inserted{(On n'a}
% L13
tjrs pas utilisé d'hyp. spéciale sur $S$\ldots).

% L14
Pour vérifier que l'on a des iso
% L15
au niveau des $H^i(B_{X,U},F)\longrightarrow H^i(X_{\mathrm{ét}},F)$,
% L16
il suffit par le lemme des cinq de
% L17
le prouver au niveau des
% L18
$H^i(B_{\pi_1,U},F)\longrightarrow H^i(U_{\mathrm{ét}},F)$ [i.e. \unsure{vérifier|vérifions} que la
% L19
cohomologie de $U_{\mathrm{ét}}$ ``se calcule galoisiennement'']
% L20
\uline{et} au niveau des $H^i_S$. \deleted{\unsure{Or}} Regardons,

\scanpage{128}
% L01
d'abord ces derniers -- on a
% L02
$H^0_S(B_{X,U},E)={}$ ens des sections de $F$ [i.e. des
% L03
couples d'une section de $F_U$ et
% L04
d'une de $F_D$, se correspondant
% L05
par $\varphi$, i.e. d'un élément \inserted{$x$} de
% L06
$E_U$ invariant par $\mathcal E$, et des
% L07
$x_i\in E_i$ invariants par les $\Gamma_i$,
% L08
tels que $x=\varphi_i(x_i)\quad\forall i$]
% L09
tels que $x=0$ \deleted{\unsure{des|dans} \unsure{systèmes}}
% L10
\[
=\prod_i(\operatorname{Ker}(E_i\longrightarrow E_U)
\]
% L11
on constate que ce foncteur se factorise
% L12
par le foncteur $H^0_S$ relatif à ``l'inclusion''
% L13
\deleted{\unsure{$D^{*}\subset D$} \unsure{ou} $B_{\pi_1,D^{*}}\longrightarrow B_{\pi_1,D}$}
% L14
analogue \deleted{\illegible{1} $B_{\pi_1,U}$ \illegible{2}}
% L15
\[
B_{\underbrace{\pi_1,\unsure{D^{*}|D_{*}}}_{\text{remplaçant }\pi_1,U}}\longrightarrow B_{D,D^{*}}\quad\ldots
\]
% L16
et les $H^i_S$ sur $D_{X,U}$ se calculent comme
% L17
\deleted{\unsure{qui n'est autre que}}
% L18
ceux dans la situation locale $B_{D\deleted{\illegible{1}},D^{*}\deleted{\illegible{1}}}$ --
% L19
où on trouve le calcul habituel en
% L20
termes de l'hom. de groupes $\mathcal E_i\longrightarrow\Gamma_i$.
\margin{que j'ai un peu oublié !}

% L21
Mais dans le cas actuel, l'hyp. $\dim_s X=1$
% L22
aux pts de $s\in S$, \deleted{et $X$ \unsure{rég.}} implique que
% L23
la situation \inserted{de topos} localisée $D^{*}_{\unsure{a|\mathrm{ét}}}\hookrightarrow D_{\unsure{a|\mathrm{ét}}}$ est déjà entièrement
% L24
définie en termes des topos fondamentaux \unsure{des}

\scanpage{129}
\mspage{63}
% L01
\deleted{\unsure{des termes de}}
% L02
i.e. le topos $B_{D,D^{*}}$ est équivalent à
% L03
$D_{\mathrm{ét}}$ -- or les $H^i_S(X_{\mathrm{ét}},F)$ se calculent
% L04
bien sur $D_{\mathrm{ét}}$\ldots On trouve donc

% L05
\uline{Prop.} L'hom. de suites exactes de coho-
% L06
mologie envisagé est un isomorphisme,
% L07
pourvu que l'on sache que l'hom.
% L08
$H^i(\mathcal E,F_{\mathcal E})\xrightarrow{\sim}H^i(U_{\mathrm{ét}},F_U)$ est un isom.
% L09
pour $\forall i$ ($F_U$ faisceau \unsure{étale|ét. loc.} sur $U_{\mathrm{ét}}$
% L10
défini par un groupe $F_{\mathcal E}$ sur lequel $\mathcal E=\pi_1(U)$
% L11
opère).

% L12
\uline{Exemple} : \inserted{OK} Si $U$ est une courbe affine
% L13
sur $k$ alg. clos \inserted{$F$ \deleted{\illegible{1}} premier à car} (car OK pour $i=0,1$, et
% L14
pour $i=2$, les deux \unsure{mbres} sont nuls
% L15
(la \unsure{ch.|coh.} galoisienne, car le groupe fondamental
% L16
\inserted{premier à $p$} est \deleted{pro} libre\ldots) -- d'où on déduit le
% L17
cas analogue pour $\unsure{U|X}$ affine sur $k$ quelcon-
% L18
que, puis = si $U$ n'est pas affine mais
% L19
simplement quasi-projective ? \deleted{\illegible{1}}
\margin{\uline{Cor.} Dans ce cas, la cohomologie de $X$ à coefficients dans un $F_X$ loc. const. \unsure{se} calcule également galoisiennement ? Non, à cause du genre $0$ !!! Dans le cas des courbes projectives lisses de genre $\ne0$ il faut un argument spécial.}

% L20
Nous nous intéressons maintenant au
% L21
cas où $X$ projective \inserted{(connexe)} sur $k$ alg. clos, \unsure{on voit donc|on sait donc}

\scanpage{130}
% L01
que : \deleted{les $H^i_!(U,F$}
% L02
a) Si $S=\varnothing$ i.e. $X=U$, la cohomologie de $U$
% L03
à (coefficients dans les syst. locaux) est celle
% L04
de $\pi$. Donc\ldots un isom. can.
% L05
\[
H^2(\pi,T)\simeq H^2(X,T)\quad(\simeq\widehat{\mathbb Z})
\]
\margin{avec \unsure{grain|grains} de sel en car. $p>0$}
% L06
\deleted{b) Si $S=$}

% L07
b) Si $S\ne\varnothing$, les $H^i_!(U,-)$ se \unsure{déterminent|décrivent}
% L08
et calculent par voie galoisienne, en termes
% L09
des \deleted{\unsure{syst.}} \inserted{groupes} à lacets, permettant de reconsti-
% L10
tuer la situation
% L11
\diagram{\begin{tikzcd}[ampersand replacement=\&,column sep=large]
 \& B_{\pi_1,D^{*}} \arrow[dl] \arrow[dr] \& \\
 B_{\pi_1,U} \arrow[dr,dotted,no head] \& \&
 {\begin{gathered}
 B_{\pi_1,D}\ (={\text{topos discret défini}}\\
 {\text{par }}S=I)
 \end{gathered}} \arrow[dl,dashed] \\
 \& B_{U,X} \&
\end{tikzcd}}
% L12
en postulant que le composant $B_{\pi_1,D_i^{*}}$
% L13
n'est autre que $B_{L_i}$ (avec les
% L14
\unsure{notations} du \unsure{n$^{o}$|cas} précédent). On
% L15
trouve alors des iso can. (par calculs locaux)
% L16
\[
H^2_{\unsure{S_i|s_i}}(B_{U,X},\unsure{L_i})\simeq\widehat{\mathbb Z}
\quad\text{i.e.}\quad
H^2_{\unsure{S_i|s_i}}(B_{U,X},\widehat{\mathbb Z})\simeq L_i^{\otimes-1}
\]
% L17
(si $g\ne0$)
\[
H^2_{\unsure{S_i|s_i}}(B_{U,X},\widehat{\mathbb Z})
\xrightarrow{\sim}H^2(X,\widehat{\mathbb Z})
\simeq H^2(\pi,\widehat{\mathbb Z})
\]
% L18
d'où \unsure{en} \deleted{\unsure{résumé}} mettant ensemble, des iso
% L19
canoniques
\[
\deleted{\unsure{H}}\quad L_i\simeq H^2(\pi,\widehat{\mathbb Z})^{\otimes-1}.
\]
\margin{Mais dans tous les cas (= si $g=0$, du moment \unsure{qu'on} n'a pas $\nu=0$) \deleted{\illegible{1}} on trouve $H^{2}_{!}(B_{\pi_1,U},\widehat{\mathbb Z})\simeq\unsure{\widehat T_{\pi}^{\otimes-1}|T_{\pi}^{\otimes-1}}$ canoniquement, d'où une description \unsure{cohom.} des \unsure{modules d'orientations}, commune au \unsure{cas sans trous et avec} trous\ldots}
